% Generated by `python -m experiments.experiment_4.appendix_captions` -- do not edit by hand.
% Methods/appendix text for fig:exp4-main (main_figure.tex). Place it under a heading
% carrying the label, e.g. \subsection{Encoder probe: details}\label{app:exp4-main-details}
% TODO(verify): Kornblith et al. (2019) CKA; Song et al. (2012) unbiased HSIC; Krajbich et al. (2010) aDDM.
\emph{Probe.} The aDDM was not in the encoder's training mix, and its attentional weight $\omega$ acts through a within-trial fixation process for which the encoder has no input channel. Gaze always starts on item~0, which turns $\omega$ into a choice bias, and with an intercept-only design and constant item values the drift $v$, drift advantage $v_{\mathrm{diff}}$ and $\omega$ are not jointly identified, so $v$ and the drift variability $s_v$ are fixed at their prior medians ($v = 1.0375$, $s_v = 0.3133$ on the natural scale), leaving $v_{\mathrm{diff}}$, $\omega$, the threshold $a$, non-decision time $\tau$ and its variability $s_\tau$. The \emph{frozen} arm takes the Experiment~1 network's encoder (trained on DDM, RDM and CDM) with its weights frozen and uses the mean of its output tokens (128 dimensions) as the summary; the \emph{baseline} arm learns a BayesFlow SetTransformer summary (16 dimensions) jointly with its posterior. Both arms read the same per-trial input, use a FlowMatching posterior at library defaults and train for 500 epochs of 100 batches of 64, i.e.\ on the same 3.2~million simulated datasets of 500 trials; they are matched on data budget and input, not on parameter count. Recovery (Figure~\ref{fig:exp4-main}, A) uses each arm's own 500 validation datasets with 500 posterior draws, and the two rows share each column's axis limits. Calibration does not separate the arms (Figures~\ref{fig:app4-arm-baseline} and~\ref{fig:app4-arm-frozen}); a paired comparison on shared datasets does (Figure~\ref{fig:app4-compare}). \emph{Geometry.} Each family is simulated with an intercept-only design and 500 trials: 200 parameter draws, each simulated 4 times independently, encoded by the original and by the LBA-fine-tuned encoder (Experiment~3, 1000 epochs on all four families) on the identical datasets. The aDDM is the same five-parameter model as in the probe, and the LBA has its start-point range fixed at $A = 0.5$. A point is one parameter draw's summary averaged over its repeats; dimensions are standardised over the pooled set of both encoders, and a single metric MDS is fitted to all points and then split by encoder. Contours are Gaussian kernel density estimates, drawn from 22\% of each family's peak. \emph{Similarity.} The Gram block uses 200 aDDM parameter draws (4 repeats each, averaged): each summary's centred Gram matrix is scaled to unit norm over its off-diagonal entries with the diagonal set to zero, and both are divided by one common constant so that they lie in $[-1, 1]$; the difference is SetTransformer minus CogFormer. CKA is linear CKA with the unbiased HSIC estimator, $\mathrm{CKA} = 0.729$. The summaries differ in dimension (128 vs 16), so for scale: random linear projections of the frozen summary itself to 16 dimensions, which keep only the frozen encoder's own geometry, score $0.795 \pm 0.056$ (30 projections; Figure~\ref{fig:app4-rsa}).
For the structure-conditioned checkpoint, this experiment probes the data-only encoder with structural context omitted. It does not evaluate the full structure-conditioned posterior encoding path.
